\chapter{从物理结构到自反结构}
\begin{chapteroverviewbox}
这一章，我们围绕“从物理结构到自反结构”做一件克制的事：把直觉压缩为状态、关系、时间尺度和可检验后果。它在全书中的任务，是说明组织、文明与自反结构演化。
\end{chapteroverviewbox}
\ChapterArt{\number\value{chapter}}

\section*{理论来路与依赖接口}
本章的直接思想来源是原文后期关于 CSO 外化、组织、文明、宇宙粗粒化与功能自反性的综合讨论。我们采用原文较晚形成的定义来整理较早讨论，不把对话中的试探性比喻与最终模型并列成互相竞争的“定律”。已有理论锚点主要是集体记忆、制度演化、分布式控制与科学认识论（Hutchins、Clark 与 Wiener）。

本章使用上一章“宇宙的两类宏观结构”已经得到的对象与边界；本章产出将供下一章“智能的最终结构定义”使用。若后文术语偶尔出现，只作为路线提示，不参与本章推导。

\section*{从哪里出发}
我们已经拥有上一章给出的概念，现在只向前走一步。我会先把对象放进闭环，再讨论它的数学形式，最后检查工程边界。读者不必预先相信本章术语；只需暂时接受一个方法论约定：智能结构的意义来自它在现实反馈中的作用，而不是来自名称本身。

令系统在观察窗口 $[t,t+T]$ 内的状态轨迹为 $x_{[t,t+T]}$，环境扰动为 $d$，行动为 $u$。本章讨论的对象若确实有效，就应改变条件分布 $p(x_{t+T}\mid x_t,u,d)$，或改变达到同一结果所需的代价。这个起点把讨论锚定在集体记忆、制度演化、分布式控制与科学认识论上。
\section{Structure Exists}
\begin{narration}
现在把镜头移到“Structure Exists”。我们只沿本章已经取得的概念向前一步：先确定它在闭环里的角色，再检查它的尺度与失败边界。
\end{narration}

本节把“Structure Exists”落实为用粗粒化和有效结构讨论认知结构在物理世界中的承载位置。这个定义不是说“任何相似现象都算”，而是要求我们同时给出对象域、允许的干预以及观察窗口。对于主题“从物理结构到自反结构”而言，它承担的是组织、文明与自反结构演化中的一个可辨识环节。

把这一关系写成最小数学骨架，可得

\begin{equation}\mathcal C_\ell=\mathsf{CoarseGrain}_\ell(\mathcal W).\end{equation}

式中的符号只承担结构角色，具体参数仍需由任务和身体数据辨识。我们首先检查可观测量，再检查控制量，最后检查比它更慢的约束。这样的顺序来自集体记忆、制度演化、分布式控制与科学认识论，但本书对Structure Exists的命名和组合仍是需要实验验证的建模提案。

这一小节的反例同样重要：从功能自反跳到宇宙主观意识属于没有证据的范畴跃迁。因此评价它不能只看一次任务结果，而要比较干预前后的轨迹、恢复时间、维护成本和风险。如果改变该机制而所有允许观察下的分布都不变，那么它至少在当前实验族中是冗余描述。

\begin{thought}
对“Structure Exists”做一次消融：冻结它、删除它，或只允许它以相邻时间尺度交换残差。哪些现象立即消失，哪些只在长期出现？若答案无法区分三种干预，我们还没有找到正确的状态变量。
\end{thought}

最后，把结论交还现实：本节只建立功能层关系，不由此推出特定脑区实现、主观意识或宇宙目的。可测状态、干预后果和明确失败条件，才是从原文直觉走向可检验理论的桥。

\section{Structure Maintains Itself}
\begin{narration}
现在把镜头移到“Structure Maintains Itself”。我们只沿本章已经取得的概念向前一步：先确定它在闭环里的角色，再检查它的尺度与失败边界。
\end{narration}

本节把“Structure Maintains Itself”落实为用粗粒化和有效结构讨论认知结构在物理世界中的承载位置。这个定义不是说“任何相似现象都算”，而是要求我们同时给出对象域、允许的干预以及观察窗口。对于主题“从物理结构到自反结构”而言，它承担的是组织、文明与自反结构演化中的一个可辨识环节。

把这一关系写成最小数学骨架，可得

\begin{equation}\mathcal C_\ell=\mathsf{CoarseGrain}_\ell(\mathcal W).\end{equation}

式中的符号只承担结构角色，具体参数仍需由任务和身体数据辨识。我们首先检查可观测量，再检查控制量，最后检查比它更慢的约束。这样的顺序来自集体记忆、制度演化、分布式控制与科学认识论，但本书对Structure Maintains Itself的命名和组合仍是需要实验验证的建模提案。

这一小节的反例同样重要：从功能自反跳到宇宙主观意识属于没有证据的范畴跃迁。因此评价它不能只看一次任务结果，而要比较干预前后的轨迹、恢复时间、维护成本和风险。如果改变该机制而所有允许观察下的分布都不变，那么它至少在当前实验族中是冗余描述。

\begin{thought}
对“Structure Maintains Itself”做一次消融：冻结它、删除它，或只允许它以相邻时间尺度交换残差。哪些现象立即消失，哪些只在长期出现？若答案无法区分三种干预，我们还没有找到正确的状态变量。
\end{thought}

最后，把结论交还现实：本节只建立功能层关系，不由此推出特定脑区实现、主观意识或宇宙目的。可测状态、干预后果和明确失败条件，才是从原文直觉走向可检验理论的桥。

\section{Structure Senses}
\begin{narration}
现在把镜头移到“Structure Senses”。我们只沿本章已经取得的概念向前一步：先确定它在闭环里的角色，再检查它的尺度与失败边界。
\end{narration}

本节把“Structure Senses”落实为用粗粒化和有效结构讨论认知结构在物理世界中的承载位置。这个定义不是说“任何相似现象都算”，而是要求我们同时给出对象域、允许的干预以及观察窗口。对于主题“从物理结构到自反结构”而言，它承担的是组织、文明与自反结构演化中的一个可辨识环节。

把这一关系写成最小数学骨架，可得

\begin{equation}\mathcal C_\ell=\mathsf{CoarseGrain}_\ell(\mathcal W).\end{equation}

式中的符号只承担结构角色，具体参数仍需由任务和身体数据辨识。我们首先检查可观测量，再检查控制量，最后检查比它更慢的约束。这样的顺序来自集体记忆、制度演化、分布式控制与科学认识论，但本书对Structure Senses的命名和组合仍是需要实验验证的建模提案。

这一小节的反例同样重要：从功能自反跳到宇宙主观意识属于没有证据的范畴跃迁。因此评价它不能只看一次任务结果，而要比较干预前后的轨迹、恢复时间、维护成本和风险。如果改变该机制而所有允许观察下的分布都不变，那么它至少在当前实验族中是冗余描述。

\begin{thought}
对“Structure Senses”做一次消融：冻结它、删除它，或只允许它以相邻时间尺度交换残差。哪些现象立即消失，哪些只在长期出现？若答案无法区分三种干预，我们还没有找到正确的状态变量。
\end{thought}

最后，把结论交还现实：本节只建立功能层关系，不由此推出特定脑区实现、主观意识或宇宙目的。可测状态、干预后果和明确失败条件，才是从原文直觉走向可检验理论的桥。

\section{Structure Represents}
\begin{narration}
现在把镜头移到“Structure Represents”。我们只沿本章已经取得的概念向前一步：先确定它在闭环里的角色，再检查它的尺度与失败边界。
\end{narration}

本节把“Structure Represents”落实为用粗粒化和有效结构讨论认知结构在物理世界中的承载位置。这个定义不是说“任何相似现象都算”，而是要求我们同时给出对象域、允许的干预以及观察窗口。对于主题“从物理结构到自反结构”而言，它承担的是组织、文明与自反结构演化中的一个可辨识环节。

把这一关系写成最小数学骨架，可得

\begin{equation}\mathcal C_\ell=\mathsf{CoarseGrain}_\ell(\mathcal W).\end{equation}

式中的符号只承担结构角色，具体参数仍需由任务和身体数据辨识。我们首先检查可观测量，再检查控制量，最后检查比它更慢的约束。这样的顺序来自集体记忆、制度演化、分布式控制与科学认识论，但本书对Structure Represents的命名和组合仍是需要实验验证的建模提案。

这一小节的反例同样重要：从功能自反跳到宇宙主观意识属于没有证据的范畴跃迁。因此评价它不能只看一次任务结果，而要比较干预前后的轨迹、恢复时间、维护成本和风险。如果改变该机制而所有允许观察下的分布都不变，那么它至少在当前实验族中是冗余描述。

\begin{thought}
对“Structure Represents”做一次消融：冻结它、删除它，或只允许它以相邻时间尺度交换残差。哪些现象立即消失，哪些只在长期出现？若答案无法区分三种干预，我们还没有找到正确的状态变量。
\end{thought}

最后，把结论交还现实：本节只建立功能层关系，不由此推出特定脑区实现、主观意识或宇宙目的。可测状态、干预后果和明确失败条件，才是从原文直觉走向可检验理论的桥。

\section{Structure Represents Itself}
\begin{narration}
现在把镜头移到“Structure Represents Itself”。我们只沿本章已经取得的概念向前一步：先确定它在闭环里的角色，再检查它的尺度与失败边界。
\end{narration}

本节把“Structure Represents Itself”落实为用粗粒化和有效结构讨论认知结构在物理世界中的承载位置。这个定义不是说“任何相似现象都算”，而是要求我们同时给出对象域、允许的干预以及观察窗口。对于主题“从物理结构到自反结构”而言，它承担的是组织、文明与自反结构演化中的一个可辨识环节。

把这一关系写成最小数学骨架，可得

\begin{equation}\mathcal C_\ell=\mathsf{CoarseGrain}_\ell(\mathcal W).\end{equation}

式中的符号只承担结构角色，具体参数仍需由任务和身体数据辨识。我们首先检查可观测量，再检查控制量，最后检查比它更慢的约束。这样的顺序来自集体记忆、制度演化、分布式控制与科学认识论，但本书对Structure Represents Itself的命名和组合仍是需要实验验证的建模提案。

这一小节的反例同样重要：从功能自反跳到宇宙主观意识属于没有证据的范畴跃迁。因此评价它不能只看一次任务结果，而要比较干预前后的轨迹、恢复时间、维护成本和风险。如果改变该机制而所有允许观察下的分布都不变，那么它至少在当前实验族中是冗余描述。

\begin{thought}
对“Structure Represents Itself”做一次消融：冻结它、删除它，或只允许它以相邻时间尺度交换残差。哪些现象立即消失，哪些只在长期出现？若答案无法区分三种干预，我们还没有找到正确的状态变量。
\end{thought}

最后，把结论交还现实：本节只建立功能层关系，不由此推出特定脑区实现、主观意识或宇宙目的。可测状态、干预后果和明确失败条件，才是从原文直觉走向可检验理论的桥。

\section{Structure Changes Reality According to Representation}
\begin{narration}
现在把镜头移到“Structure Changes Reality According to Representation”。我们只沿本章已经取得的概念向前一步：先确定它在闭环里的角色，再检查它的尺度与失败边界。
\end{narration}

本节把“Structure Changes Reality According to Representation”落实为用粗粒化和有效结构讨论认知结构在物理世界中的承载位置。这个定义不是说“任何相似现象都算”，而是要求我们同时给出对象域、允许的干预以及观察窗口。对于主题“从物理结构到自反结构”而言，它承担的是组织、文明与自反结构演化中的一个可辨识环节。

把这一关系写成最小数学骨架，可得

\begin{equation}\mathcal C_\ell=\mathsf{CoarseGrain}_\ell(\mathcal W).\end{equation}

式中的符号只承担结构角色，具体参数仍需由任务和身体数据辨识。我们首先检查可观测量，再检查控制量，最后检查比它更慢的约束。这样的顺序来自集体记忆、制度演化、分布式控制与科学认识论，但本书对Structure Changes Reality According to Representation的命名和组合仍是需要实验验证的建模提案。

这一小节的反例同样重要：从功能自反跳到宇宙主观意识属于没有证据的范畴跃迁。因此评价它不能只看一次任务结果，而要比较干预前后的轨迹、恢复时间、维护成本和风险。如果改变该机制而所有允许观察下的分布都不变，那么它至少在当前实验族中是冗余描述。

\begin{thought}
对“Structure Changes Reality According to Representation”做一次消融：冻结它、删除它，或只允许它以相邻时间尺度交换残差。哪些现象立即消失，哪些只在长期出现？若答案无法区分三种干预，我们还没有找到正确的状态变量。
\end{thought}

最后，把结论交还现实：本节只建立功能层关系，不由此推出特定脑区实现、主观意识或宇宙目的。可测状态、干预后果和明确失败条件，才是从原文直觉走向可检验理论的桥。

\section{Structure Changes How It Changes Itself}
\begin{narration}
现在把镜头移到“Structure Changes How It Changes Itself”。我们只沿本章已经取得的概念向前一步：先确定它在闭环里的角色，再检查它的尺度与失败边界。
\end{narration}

本节把“Structure Changes How It Changes Itself”落实为用粗粒化和有效结构讨论认知结构在物理世界中的承载位置。这个定义不是说“任何相似现象都算”，而是要求我们同时给出对象域、允许的干预以及观察窗口。对于主题“从物理结构到自反结构”而言，它承担的是组织、文明与自反结构演化中的一个可辨识环节。

把这一关系写成最小数学骨架，可得

\begin{equation}\mathcal C_\ell=\mathsf{CoarseGrain}_\ell(\mathcal W).\end{equation}

式中的符号只承担结构角色，具体参数仍需由任务和身体数据辨识。我们首先检查可观测量，再检查控制量，最后检查比它更慢的约束。这样的顺序来自集体记忆、制度演化、分布式控制与科学认识论，但本书对Structure Changes How It Changes Itself的命名和组合仍是需要实验验证的建模提案。

这一小节的反例同样重要：从功能自反跳到宇宙主观意识属于没有证据的范畴跃迁。因此评价它不能只看一次任务结果，而要比较干预前后的轨迹、恢复时间、维护成本和风险。如果改变该机制而所有允许观察下的分布都不变，那么它至少在当前实验族中是冗余描述。

\begin{thought}
对“Structure Changes How It Changes Itself”做一次消融：冻结它、删除它，或只允许它以相邻时间尺度交换残差。哪些现象立即消失，哪些只在长期出现？若答案无法区分三种干预，我们还没有找到正确的状态变量。
\end{thought}

最后，把结论交还现实：本节只建立功能层关系，不由此推出特定脑区实现、主观意识或宇宙目的。可测状态、干预后果和明确失败条件，才是从原文直觉走向可检验理论的桥。

\section*{本章收束：把名词还原为闭环}
现在回头看“从物理结构到自反结构”，最重要的不是记住缩写，而是说清本章对象怎样改变系统的状态、代价或可恢复性。下一章“智能的最终结构定义”只使用这里已经给出的结果，不反过来为本章提供前提。

\begin{bookproposition}
若一个候选机制既不改变可观测轨迹分布，也不改变资源、风险或恢复代价，那么在当前观察尺度上，它不是一个可辨识的功能结构。
\end{bookproposition}
\begin{proofidea}
功能等价由可干预后果定义。若所有允许干预下的输出分布与代价均不变，则该机制相对于当前实验族不可辨识；要么扩大观察尺度，要么删除这一冗余描述。
\end{proofidea}

\begin{readercheck}
请尝试回答：本章对象的状态是什么？它的最快和最慢时间常数是什么？什么观测能证伪它？它失败时，残差应向哪一层升级？
\end{readercheck}
\cleardoublepage
